%!TEX root = ../main.tex
% =============================================================================
%  五、问题 2：不确定性下的两阶段计划购电模型     建议 4–5 页
%  评审在问：不确定性处理合理吗？
%  本节必须出现：① 信息结构声明（R3：附件 3 的 0:00 预报 + 负载已知；结算用附件 2 实际值）
%                ② 两阶段模型全式（第一阶段 E/c/d 与场景无关；第二阶段追索 g）
%                ③ 场景生成与削减三层（残差定义 / 按日总残差分层 / 层内 medoid / 层权重）
%                ④ 报童型安全边际的推导、两界、四条适用边界与实测覆盖率
%                ⑤ 求解结果：计划费 / 紧急费 / 合计 / 紧急电量与触发天数
%                ⑥ 四个指定日期的表 1/表 2/表 3
%                ⑦ 对照：完美信息 R1 与典型日预报 R2
%                ⑧ 紧急购电的时标解释（发生在结算时刻，不在决策时刻）
%  口径与图名以 notes/写作交底.md §六 为准：C-05 在本节首次定义；图名为 fig07_q2_emergency.pdf。
%  红线：报童模型只作"近似刻画 + 上下界 + 方向性解释"，不得声称它是精确解。
% =============================================================================

\section{问题 2：不确定性下的两阶段计划购电模型}

\subsection{问题分析与决策结构}

问题 2 与问题 1 的差别有三点：电价仍逐日相同（附件 1 固定价，均值 0.7662~\unit{元/kWh}），
但负载与光伏\textbf{逐日逐时段变化}；计划仍在每天 \(0{:}00\) 制定；若微网提供的电能低于负载，
须按\textbf{交易时刻电价的 5 倍}紧急购电。故决策分两个层次：

\begin{itemize}
  \item \textbf{第一阶段（日前 \(0{:}00\)）}：只知预报，一次性确定全天 144 个时段的
        \textbf{计划购电量 \(E_t\)}与\textbf{充放电量 \(c_t,d_t\)}，该决策\textbf{不依赖场景}；
  \item \textbf{第二阶段（实时结算）}：实际值揭示后，计划与实际之差由
        \textbf{紧急购电量 \(g_{t,\omega}\)}（缺口，按 \(5p_t\) 计价）与
        \textbf{光伏盈余及超购部分的处置}吸收。
\end{itemize}

\paragraph{信息结构声明}
本文采用\textbf{信息结构 R3}：计划依据为\textbf{附件 3 中当日 \(0{:}00\) 发布的光伏功率预报}
\(\widehat{PV}^{\text{fcst}}_t=\widehat{PV}^{(0{:}00)}_{t}\)；题面未给负载预报，故计划阶段按
"当日负载曲线已知"处理（等价于负载预报无误差），并在假设节显式声明；\textbf{结算依据为附件 2
的当日实际负载与实际光伏}。两者是\textbf{两套数据}，其差异即紧急购电的来源。作为对照，另设完美信息 R1
（\(0{:}00\) 即知当日真值）与典型日预报 R2（附件 1 的日均形状），结果见 \S\ref{sec:q2-compare}。

\paragraph{时标解释：紧急购电发生在结算时刻}
紧急购电不是在某个决策时刻"决定"的，而是在\textbf{结算时刻}由缺口
\(g_{t,\omega}=\max\{0,(L_{\omega,t}-PV_{\omega,t})\Delta t+c_t-d_t-E_t\}\) 唯一确定的：
计划量照常执行，缺口只能以 \(5p_t\) 的价格即时补足，不能回溯改写计划。
这一时标差异是下文报童模型第三条适用边界的来源（\S\ref{sec:q2-newsvendor}）。

\paragraph{一个必须显式声明的建模选择}
\textbf{充放电计划同样在第一阶段确定，实时不允许重调度}。依据是题面"其他时间段的购电费用均按
计划购电量计算"，以及结果文件要求填报逐日逐时段的充放电量。据此，第一阶段买到的高价电不能因
实际光伏更好而转手，计划购电量\textbf{允许高于}预测口径的需求：超出部分既不产生售电收益、
也不再被后续时段使用（超购即被处置），其成本为 \(p_t\)。这一"超购即弃"的设定使模型天然具有
\textbf{安全边际}，也是报童型权衡在本问成立的结构性原因（经济价值见 \S\ref{sec:q2-newsvendor}）。

\subsection{场景生成与削减}

光伏预报误差是本问唯一的不确定性来源。定义\textbf{逐日逐时段残差}
\begin{equation}
  r_{d,t}=\bigl(\widehat{PV}^{\text{fcst}}_{d,t}-PV^{\text{act}}_{d,t}\bigr)\Delta t,
  \qquad d\in\mathcal D,\ t\in\mathcal T ,
  \label{eq:q2-resid}
\end{equation}
单位为 \unit{kWh}，样本为窗口内全部 \(|\mathcal D|=334\) 天 \(\times\,144\) 槽。场景集分三层生成：

\begin{enumerate}
  \item \textbf{按日总残差分层}：按 \(s_d=\sum_{t}r_{d,t}\) 升序把 334 天\textbf{等频}划分为
        \(K\) 层（主口径 \(K=60\)，每层 5--6 天），保留残差水平分布的\textbf{全部区间}（尾部不被均值化抹平）；
  \item \textbf{层内取 medoid}：每层取"到本层其他样本欧氏距离之和最小"的
        \textbf{真实历史日残差向量}（不做中心化与平滑），逐槽保留残差的\textbf{日内形状}
        与时段间相关结构；
  \item \textbf{权重 = 层样本占比}：\(W_\omega=n_\omega/334\)，\(\sum_\omega W_\omega=1\)。
\end{enumerate}

情景 \(\omega\) 下实际净需求 \(ND_{\omega,t}=ND^{\text{fcst}}_t+r_{\omega,t}\)，
其中 \(ND^{\text{fcst}}_t=(L^{\text{fcst}}_t-\widehat{PV}^{\text{fcst}}_t)\Delta t\)。
场景集一次性生成后\textbf{存盘固定}，问题 3、4 与全部消融实验读取同一份，以保证"改进"来自
模型而非场景噪声。



三种层数下总费用极差 13.65 万元，仅为主口径的 \(0.96\%\)，\textbf{结论对场景层数稳健}；
紧急购电量对 \(K\) 更敏感（\(K=40\) 时为 \(K=120\) 时的 1.68 倍）：层数越少、单层样本越多，代表场景越难贴近尾部。场景集退化为"残差恒为零"的单一场景
（\(K=1\)，即严格按预报口径购电）即下文的\textbf{无安全边际}消融（\S\ref{sec:q2-newsvendor}）。

\subsection{两阶段随机规划模型}

\begin{subequations}\label{eq:q2}
  % 目标：第一阶段购电费 + 第二阶段紧急购电费的期望
  \begin{equation}
    \min_{E,c,d,S}\quad
    \sum_{t=1}^{144}p_t E_t
    \;+\;\sum_{\omega\in\Omega}W_\omega\sum_{t=1}^{144}5\,p_t\,g_{t,\omega}
    \label{eq:q2-obj}
  \end{equation}
  % 第一阶段决策（非预期性）
  \begin{equation}
    E_t,\ c_t,\ d_t,\ S_t\ \text{与}\ \omega\ \text{无关},
    \qquad \forall t\in\mathcal T
    \label{eq:q2-nonanticip}
  \end{equation}
  % 计划可行性（允许超购）
  \begin{equation}
    E_t-c_t+d_t\ \ge\ \bigl(L^{\text{fcst}}_t-\widehat{PV}^{\text{fcst}}_t\bigr)\Delta t,
    \qquad \forall t\in\mathcal T
    \label{eq:q2-plan}
  \end{equation}
  % 第二阶段：缺口由紧急购电补齐
  \begin{equation}
    g_{t,\omega}\ \ge\ \bigl(L_{\omega,t}-PV_{\omega,t}\bigr)\Delta t+c_t-d_t-E_t,
    \qquad g_{t,\omega}\ge 0
    \label{eq:q2-gap}
  \end{equation}
  % 储电量递推、界与终端条件
  \begin{equation}
    S_t=S_{t-1}+\eta_c c_t-\frac{d_t}{\eta_d},\quad
    \underline S\le S_t\le\overline S,\quad
    0\le c_t,d_t\le\overline P\Delta t,\quad
    S_0=S^{init},\quad S_{144}\ge S_0
    \label{eq:q2-soc}
  \end{equation}
  % 非负
  \begin{equation}
    E_t\ge 0,\qquad g_{t,\omega}\ge 0
    \label{eq:q2-nonneg}
  \end{equation}
\end{subequations}

式~\eqref{eq:q2-nonanticip} 是两阶段结构的代数表达：带 \(\omega\) 下标的量为\textbf{事后}量
（结算量），第一阶段变量 \(E_t,c_t,d_t,S_t\) 在全部 60 个场景间共享。式~\eqref{eq:q2-plan} 用
\textbf{不等号}：计划购电量只要求不低于"预测口径净需求 + 充电量 − 放电量"；式~\eqref{eq:q2-gap}
的缺口由紧急购电补齐，\(5p_t>0\) 使最优解中 \(g_{t,\omega}\) 取其下界，即
\begin{equation}
  g_{t,\omega}=\max\Bigl\{0,\ \bigl(L_{\omega,t}-PV_{\omega,t}\bigr)\Delta t+c_t-d_t-E_t\Bigr\},
  \label{eq:q2-gapmax}
\end{equation}
是式~\eqref{eq:q2-gap} 的紧形式；式~\eqref{eq:q2-gap} 用的是\textbf{计划充放电量}
\(c_t,d_t\)（不含 \(\omega\)），与 \S5.1 一致。模型规模（\(K=60\)）为每个子问题 \(576+144\times60=9216\)
个变量、144 个等式与 \(144\times(1+60)\) 个不等式，仍是纯 LP，用 SciPy 的 HiGHS 求解器精确求解；
结果窗口内共求解 334 个子问题（逐日顺序推进，日初储电量取前一日末值，1 月作预热期，假设 A8），
终端条件取 \(S_{144}\ge S_0\)（假设 A6）。

\subsection{报童型安全边际的可解释刻画}\label{sec:q2-newsvendor}

暂不考虑储能的跨期耦合，把每个时段看作独立的"报童问题"：计划购电量 \(E_t\) 相当于订货量，
当日实际需求 \(D_t=(L_t-PV_t)\Delta t+c_t-d_t\) 相当于需求量。缺 1~\unit{kWh} 要多付
\((5-1)p_t=4p_t\)（\textbf{缺货成本} \(c_u=4p_t\)，正常电价 \(p_t\) 已计入计划口径）；
多买 1~\unit{kWh} 的损失记为 \(c_o\)。由报童模型，最优计划量是需求分布的\textbf{临界分位数}：
\begin{equation}
  q^{*}_t=F_{D_t}^{-1}\!\left(\frac{c_u}{c_u+c_o}\right)
  =F_{D_t}^{-1}\!\left(\frac{4p_t}{4p_t+c_o}\right),
  \label{eq:q2-newsvendor}
\end{equation}
与之互补的是\textbf{最优缺货概率} \(1-q^{*}_t\)；两者不可混写。

过剩成本 \(c_o\) 的两种口径构成临界分位数的\textbf{上下界}：\textbf{口径一}（超购全废弃）
\(c_o=p_t\)，临界分位数 \(4/5=0.800\)、缺货概率 \(1/5=0.200\)；\textbf{口径二}（超购被储能吸纳
并在后续时段替代购电，单位回收价值 \(p^{low}\)）\(c_o=p_t-p^{low}\approx0.37\)~\unit{元/kWh}，
取全天电价均值 \(p_t=0.7662\)~\unit{元/kWh} 得临界分位数 \(3.0648/(3.0648+0.37)=0.892\)、
缺货概率 \(0.108\approx1/10\)。故\textbf{缺货概率的两个界是 \(1/10\) 与 \(1/5\)}，
覆盖率对应 \(0.892\) 与 \(0.800\)。

\begin{table}[htbp]
  \centering
  \caption{报童近似的两条界与两阶段模型的实测缺货频率（窗口 2025-02-01 至 2025-12-31）}
  \label{tab:newsvendor}
  \footnotesize
  \begin{tabularx}{\linewidth}{L{3.4cm} C{2.0cm} C{2.4cm} C{1.7cm} Y}
    \toprule
    过剩成本口径 & 临界分位数 & 最优缺货概率 & 样本量 & 与实测的关系 \\
    \midrule
    口径一：超购全废弃 \(c_o=p_t\) & 0.800 & 0.200（\(1/5\)） & --- & 缺货概率上界 \\
    口径二：超购可回收 \(c_o\approx0.37\)~\unit{元/kWh} & 0.892 & 0.108（\(\approx1/10\)） & --- & 缺货概率下界 \\
    \midrule
    两阶段随机规划实测 & 覆盖率 0.9145 & \textbf{0.0855} & 48\,096 槽 & 低于下界（更保守） \\
    \bottomrule
  \end{tabularx}
  \\[1ex]
  {\footnotesize 表注：实测缺货频率 = 紧急购电量大于 \(10^{-6}\)~\unit{kWh} 的槽数占 334 天 \(\times\) 144 槽的比例；
  覆盖率 = \(1-\) 缺货频率；电量满足率为 \(99.04\%\)（口径与 \calref{C-05} 一致）。}
\end{table}

两阶段随机规划的实测行为：\textbf{逐槽缺货频率 8.55\%}（覆盖率 91.45\%）低于下界 \(1/10\)，
\textbf{电量满足率 99.04\%}，全窗口\textbf{安全边际合计 576\,299.67~kWh}（平均 11.98~\unit{kWh}/槽）
\evid{C-Q2-01}。实测更保守的原因在结构而非参数：\textbf{储能同时承担对冲功能}，同一份计划可由
电池放电在缺口时段补足，单期报童模型无法刻画"计划购电量"与"电池吞吐"的替代关系。
故式~\eqref{eq:q2-newsvendor} 与表~\ref{tab:newsvendor} 只作\textbf{方向性近似与上下界}，
\textbf{不是}精确最优解。

安全边际的经济价值由\textbf{无安全边际消融}量化：场景集退化为单一情景（\(K=1\)，残差恒为零，
严格按预报口径购电）后，全年计划购电费 12\,570\,533.91~\unit{元}、紧急购电费
4\,828\,025.76~\unit{元}、合计 17\,398\,559.68~\unit{元}，比无储能基线 16\,407\,319.63~\unit{元}
高 \(6.04\%\)，紧急购电量 1\,156\,618.02~\unit{kWh}、334/334 天全部触发；主口径费用
14\,264\,792.13~\unit{元}、紧急购电量 172\,015.82~\unit{kWh}。分解为：\textbf{安全边际把计划购电费
推高 1\,067\,018.99~\unit{元}（边际电量 576\,299.67~\unit{kWh} 按均价 0.6176~\unit{元/kWh} 计
355\,918.38~\unit{元}，其余 711\,100.60~\unit{元} 来自储能排产与购电时序的同步调整），
却把紧急购电费压低 4\,200\,786.52~\unit{元}，净收益 3\,133\,767.55~\unit{元}
（相对不设边际降低 \(18.0\%\)）}。即边际电量只占全窗口计划购电量 22\,220\,333.14~\unit{kWh} 的
\(2.59\%\)，却消掉 \(87.0\%\) 的紧急购电费。惩罚倍数灵敏度方向一致：同一典型日与同一残差场景集上
把倍数由 \(k=2\) 提高到 \(k=10\)，单日最优安全边际由 671~\unit{kWh} 增至 3\,334~\unit{kWh}，
期望紧急购电量由 1\,246~\unit{kWh} 降至 187~\unit{kWh}，目标值由 39\,224.54~\unit{元} 升至
41\,758.99~\unit{元}（\(k=5\) 时为 40\,747.45~\unit{元}）。

报童模型在本问题中\textbf{不是精确解}，其偏离来源有四条，也就是它的适用边界：

\begin{enumerate}
  \item \textbf{多时段耦合}：144 个时段通过储电量递推耦合，本时段多买的电会占用电池容量、
        影响其后所有时段，不能拆成 144 个独立报童问题；
  \item \textbf{储能缓冲}：过剩电量可被储能吸纳并在后续高价时段替代购电，故 \(c_o\) 严格小于
        \(p_t\)，临界分位数应高于口径一的 0.800；
  \item \textbf{紧急购电的时标在结算时刻}：缺口只能事后按 \(5p_t\) 补足，受当日计划充放电量与
        放电功率限制，不能回溯改写计划；
  \item \textbf{残差分布未知}：本文用 334 天样本的\textbf{样本分位数}替代总体分位数，
        每个时段仅 334 个观测，分位数估计带误差，须报置信区间。
\end{enumerate}



\subsection{对照方案与风险维度}\label{sec:q2-compare}

在同一批实际轨迹上比较三种计划依据：

\begin{itemize}
  \item \textbf{R1 完美信息}：\(0{:}00\) 即知当日实际负载与实际光伏，紧急购电量恒为 0。
        在与主方案相同的日决策结构下逐日求解，总费用 \textbf{13\,411\,560.90~\unit{元}}，
        相对无储能基线节省 \(18.26\%\)，是"预测误差成本"的\textbf{下界}；若放弃日决策解耦、
        把 334 天写成单一 LP（单一储电量轨迹），费用 12\,228\,166.56~\unit{元}，为更松的下界参考；
  \item \textbf{R3 本文主方案}（附件 3 的 \(0{:}00\) 预报）：总费用 14\,264\,792.13~\unit{元}，
        相对基线节省 \(13.06\%\)；
  \item \textbf{R2 典型日预报}：用附件 1 的日均形状作预报，总费用
        \textbf{23\,484\,646.63~\unit{元}}，比无储能基线高 \(43.14\%\)，紧急购电量
        2\,525\,548.76~\unit{kWh}（286/334 天触发）。附件 1 的数值是附件 2 实际功率的
        \textbf{全年逐时段均值}而非日前预报，只用于说明"预报质量决定成本"，\textbf{不可作为主方案}。
\end{itemize}

R3 与 R1 之差 \(14\,264\,792.13-13\,411\,560.90=853\,231.23\)~\unit{元}（相对 R1 高 \(6.36\%\)）
即\textbf{预测误差的成本}，即"5 倍紧急购电"机制对应的金额。

风险维度上，期望费用最小化对极端场景不敏感：在 334 天中每 6 天取 1 天的子样本（56 天）上做
CVaR（置信水平 \(\alpha=0.95\)，\calref{C-11}）与 \(\lambda\) 扫描，\(\lambda\) 由 0 增至 1 时
期望费用由 41\,995.59~\unit{元}/日 升至 42\,411.56~\unit{元}/日（\(+0.99\%\)），CVaR\(_{0.95}\)
由 52\,176.36~\unit{元}/日 降至 50\,510.46~\unit{元}/日（\(-3.19\%\)），日均紧急购电量由
480.2~\unit{kWh} 降至 313.8~\unit{kWh}；完整前沿与回测见 \S7.4，本节只保留与日前决策直接相关的
R1/R2/R3 对照。

\subsection{求解结果}

在 2025-02-01 至 2025-12-31（334 天）窗口内求解式~\eqref{eq:q2}，总费用
\textbf{14\,264\,792.13~\unit{元}}（1\,426.48 万元），其中\textbf{计划购电费
13\,637\,552.90~\unit{元}、紧急购电费 627\,239.24~\unit{元}}；相对无储能基线
\(C^{base}=16\,407\,319.63\)~\unit{元}（\calref{C-01} 变体 V2）节省 2\,142\,527.50~\unit{元}，
节省率 \(13.06\%\)（\calref{C-04}），两项费用占比 \(95.60\%\) 与 \(4.40\%\)（\calref{C-03}）。
全窗口计划购电量 22\,220\,333.14~\unit{kWh}（\calref{C-02}），紧急购电量 172\,015.82~\unit{kWh}，
出现在 296/334 天（\calref{C-05}）；逐槽缺货频率 8.55\%、电量满足率 99.04\%、安全边际合计
576\,299.67~\unit{kWh}，与 \S\ref{sec:q2-newsvendor} 一致。

\begin{caliber}{C-05}{紧急购电量 \(G\)}
\textbf{定义}\quad \(G=\sum_{d\in\mathcal D}\sum_{t=1}^{144}g_{d,t}\)，
其中 \(g_{d,t}\) 由式~\eqref{eq:q2-gapmax} 在\textbf{结算时刻}确定，
是"含电池吞吐项"的完整缺口，而不是 \(\max\{0,\text{实际净负荷}-\text{计划购电量}\}\)；%
\textbf{数据}\quad 结算依据为附件 2 的实际负载与实际光伏，
计划量与充放电量取 \texttt{result2.xlsx} 的对应工作表；%
\textbf{时段}\quad 334 天 \(\times\) 144 槽（2025-02-01 至 2025-12-31）；%
\textbf{单位}\quad \unit{kWh}（表内保留 2 位小数）；%
\textbf{不含}\quad 不含问题 3 调整机制引入的缺口（调整量单独统计，\calref{C-09}）；%
\textbf{对照}\quad 完美信息 R1 下为 0.0000~\unit{kWh}（0/334 天）；
典型日预报 R2 下为 2\,525\,548.76~\unit{kWh}（286/334 天）；%
\textbf{复算}\quad \texttt{code/lp\_engine.py::realized\_slot}，审计项 A1 与 A9。
\end{caliber}

\paragraph{可行性复核}
逐日检查：计划口径残差（供给缺口）最大 \(3.41\times10^{-13}\)~\unit{kWh}，储电量净递推残差最大
\(1.09\times10^{-11}\)~\unit{kWh}，最差单日综合残差 \(1.82\times10^{-12}\)；同时充放、SOC 越界、
功率越界、负购电的时段数均为 0，末值低于初值的天数 0、费用相对复算误差 0.0000；
结果文件回读校验的全天购电量最大偏差 0.0011~\unit{kWh}、全天购电费最大相对偏差
\(3.40\times10^{-7}\)。

\paragraph{题面指定日期的结果（表 1 / 表 2 / 表 3 格式）}
四个指定日期 2025-03-20、2025-06-21、2025-09-23、2025-12-21 的计划购电量
（表~\ref{tab:q2-buy}）、储能充放电量与 \(0{:}00\)/\(24{:}00\) 储电量（表~\ref{tab:q2-soc}）、
紧急购电量（表~\ref{tab:q2-emergency}）见表。四日均满足计划口径供需平衡
（最大残差 \(1.82\times10^{-12}\)）与日周期终端条件，\(0{:}00\) 与 \(24{:}00\) 储电量均为
6000.0000~\unit{kWh}。

\begin{table}[htbp]
  \centering
  \caption{微网在指定时段与指定日期的购电量、全天购电量和购电费（问题 2；电量单位：\unit{kWh}，费用单位：\unit{元}）}
  \label{tab:q2-buy}
  \small
  \begin{tabular}{lrrrr}
    \toprule
    时段／指标 & 2025-03-20 & 2025-06-21 & 2025-09-23 & 2025-12-21 \\
    \midrule
    \(10{:}00\)--\(10{:}10\) & 0.0000 & 0.0000 & 0.0000 & 0.0000 \\
    \(12{:}00\)--\(12{:}10\) & 804.3029 & 0.0000 & 536.2415 & 997.0335 \\
    \(14{:}00\)--\(14{:}10\) & 0.0000 & 0.0000 & 0.0000 & 0.0000 \\
    \(16{:}00\)--\(16{:}10\) & 817.3061 & 255.2254 & 902.7969 & 1001.2240 \\
    \(18{:}00\)--\(18{:}10\) & 697.9973 & 0.0000 & 764.3298 & 715.4144 \\
    \(20{:}00\)--\(20{:}10\) & 0.0000 & 0.0000 & 0.0000 & 0.0000 \\
    \midrule
    全天购电量 & 74\,178.1298 & 35\,701.5391 & 74\,744.7963 & 99\,692.6601 \\
    全天购电费 & 45\,229.90 & 20\,324.26 & 47\,063.39 & 64\,376.98 \\
    紧急购电量 & 1025.8805 & 187.0843 & 193.3669 & 22.5250 \\
    全天总费用 & 49\,654.02 & 21\,255.45 & 47\,717.70 & 64\,432.28 \\
    \bottomrule
  \end{tabular}
  \\[1ex]
  {\footnotesize 表注：时段按题面表述取区间左端时刻；
  全天购电费为计划口径（\(\sum_t p_tE_t\)），全天总费用为计划购电费与紧急购电费之和；
  四个日期的无储能基线分别为 52\,290.65、29\,668.18、54\,741.30、71\,176.92~\unit{元}。}
\end{table}

\begin{table}[htbp]
  \centering
  \caption{储能设备在指定日期的充放电量及 \(0{:}00\) 与 \(24{:}00\) 的储电量（问题 2，单位：\unit{kWh}）}
  \label{tab:q2-soc}
  \footnotesize
  \begin{tabular}{L{2.0cm}L{2.1cm}rr@{\quad}L{2.0cm}L{2.1cm}rr}
    \toprule
    日期 & 时段 & 充电量 & 放电量 & 日期 & 时段 & 充电量 & 放电量 \\
    \midrule
    2025-03-20 & \(0{:}00\)--\(4{:}00\)   & 5333.3333 & 828.0303  & 2025-09-23 & \(0{:}00\)--\(4{:}00\)   & 5333.3333 & 790.6233 \\
    2025-03-20 & \(4{:}00\)--\(8{:}00\)   & 6500.0000 & 7811.9697 & 2025-09-23 & \(4{:}00\)--\(8{:}00\)   & 5833.3333 & 7778.1258 \\
    2025-03-20 & \(8{:}00\)--\(12{:}00\)  & 6905.5931 & 2218.5304 & 2025-09-23 & \(8{:}00\)--\(12{:}00\)  & 6728.4081 & 1606.2615 \\
    2025-03-20 & \(12{:}00\)--\(16{:}00\) & 5333.3333 & 8640.0000 & 2025-09-23 & \(12{:}00\)--\(16{:}00\) & 5333.3333 & 8640.0000 \\
    2025-03-20 & \(16{:}00\)--\(20{:}00\) & 0.0000    & 0.0000    & 2025-09-23 & \(16{:}00\)--\(20{:}00\) & 0.0000    & 0.0000 \\
    2025-03-20 & \(20{:}00\)--\(24{:}00\) & 0.0000    & 0.0000    & 2025-09-23 & \(20{:}00\)--\(24{:}00\) & 0.0000    & 0.0000 \\
    \midrule
    2025-06-21 & \(0{:}00\)--\(4{:}00\)   & 0.0000    & 2723.8246 & 2025-12-21 & \(0{:}00\)--\(4{:}00\)   & 5333.3333 & 0.0000 \\
    2025-06-21 & \(4{:}00\)--\(8{:}00\)   & 5348.6522 & 1710.2113 & 2025-12-21 & \(4{:}00\)--\(8{:}00\)   & 6666.6667 & 9315.0000 \\
    2025-06-21 & \(8{:}00\)--\(12{:}00\)  & 5458.7995 & 0.0000    & 2025-12-21 & \(8{:}00\)--\(12{:}00\)  & 8333.3333 & 2835.0000 \\
    2025-06-21 & \(12{:}00\)--\(16{:}00\) & 5333.3333 & 8640.0000 & 2025-12-21 & \(12{:}00\)--\(16{:}00\) & 5333.3333 & 8640.0000 \\
    2025-06-21 & \(16{:}00\)--\(20{:}00\) & 0.0000    & 0.0000    & 2025-12-21 & \(16{:}00\)--\(20{:}00\) & 0.0000    & 0.0000 \\
    2025-06-21 & \(20{:}00\)--\(24{:}00\) & 0.0000    & 0.0000    & 2025-12-21 & \(20{:}00\)--\(24{:}00\) & 0.0000    & 0.0000 \\
    \bottomrule
  \end{tabular}
  \\[1ex]
  {\footnotesize 表注：四个指定日期的 \(0{:}00\) 与 \(24{:}00\) 储电量均为 6000.0000~\unit{kWh}；
  每段含 36 个 10~\unit{min} 槽，充放电量为交流侧电量（\calref{C-07}）；
  各日充放电量之比均为 1.2346（对应往返效率 0.81）。}
\end{table}

\begin{table}[htbp]
  \centering
  \caption{微网在指定日期的紧急购电量（问题 2；购电量单位：\unit{kWh}，电价单位：\unit{元/kWh}）}
  \label{tab:q2-emergency}
  \footnotesize
  \begin{tabular}{llrr}
    \toprule
    日期 & 购电时间段 & 紧急购电量 & 交易时刻电价（加权平均） \\
    \midrule
    2025-03-20 & \(6{:}00\)--\(8{:}00\)   & 228.9914 & 0.8625 \\
    2025-03-20 & \(8{:}10\)--\(8{:}30\)   & 49.1586  & 0.8625 \\
    2025-03-20 & \(8{:}50\)--\(10{:}30\)  & 531.2739 & 0.8625 \\
    2025-03-20 & \(10{:}40\)--\(11{:}40\) & 216.4567 & 0.8625 \\
    \midrule
    2025-06-21 & \(4{:}00\)--\(4{:}50\)   & 5.2416   & 0.9955 \\
    2025-06-21 & \(14{:}00\)--\(14{:}30\) & 181.8427 & 0.9955 \\
    \midrule
    2025-09-23 & \(6{:}50\)--\(7{:}40\)   & 132.5948 & 0.6768 \\
    2025-09-23 & \(11{:}00\)--\(11{:}40\) & 60.7721  & 0.6768 \\
    \midrule
    2025-12-21 & \(11{:}00\)--\(11{:}30\) & 22.5250 & 0.4910 \\
    \midrule
    \multicolumn{2}{l}{四个日期合计} & 1428.8567 & --- \\
    \bottomrule
  \end{tabular}
  \\[1ex]
  {\footnotesize 表注：本表按"日期 \(\times\) 紧急购电事件"逐行展开，与题面表 3 的横向版式一一对应；
  同一日期可有多行，连续触发的相邻槽按事件聚合，时间段为该事件涉及的槽位区间；
  交易时刻电价为该事件内按电量加权的平均电价（事件电价按 \(5p_t\) 计价，本列为 \(p_t\)）。
  无紧急购电的时段不列入本表（全窗口 334 天中 296 天有触发）。}
\end{table}
